
\section{A Powerful Image Under Examination}

One sometimes hears a striking explanation of the Eucharist: at the Last
Supper Christ carried the apostles forward to Calvary, whereas at every later
Mass he carries the Church backward to Calvary. The image tries to defend
three Catholic convictions at once. The Supper is truly sacrificial; the Cross
is the one unrepeatable sacrifice; and the Mass is more than a lesson about an
absent event.

Those convictions are sound. The proposed temporal mechanism is not. Scripture,
the Fathers, the scholastic synthesis, and the Church's authoritative teaching
speak of institution, memorial, sacramental representation, participation,
application, and presence. They do not teach that the apostles left the upper
room's time, or that worshippers at Mass leave theirs. The image can be saved
only as consciously figurative shorthand. Taken literally, it obscures the
very sacramental realism it was meant to protect.

The better claim is both more exact and more profound. At the Supper,
before the bloody immolation on the Cross, Christ sacramentally gave under
separate signs the Body he was about to hand over and the Blood he was about
to shed. On Calvary he consummated that self-offering in history once for all.
At every Mass after Calvary, the same living Christ makes that one sacrifice
present under sacramental signs and joins the Church's offering to his. The
relations ``before'' and ``after'' remain real. The sacrifice is one because
the victim and offerer are the same: Christ makes his once-for-all
self-offering sacramentally present and offers through the ministry of priests,
while only the manner of offering differs. Its unity does not require several
congregations to occupy one clock-time.

